\documentclass[10pt,a4paper]{article}
\usepackage[T1]{fontenc}
\usepackage[utf8]{inputenc}
\usepackage{lmodern}
\usepackage[margin=2.1cm,top=1.8cm]{geometry}
\usepackage{xcolor}
\usepackage{titlesec}
\usepackage{enumitem}
\usepackage{tabularx}
\usepackage[hidelinks]{hyperref}

\definecolor{navy}{HTML}{0B1F3A}
\definecolor{gold}{HTML}{A6863F}
\pagestyle{empty}
\setlength{\parindent}{0pt}

\titleformat{\section}{\color{navy}\normalsize\bfseries}{}{0pt}{}[{\color{gold}\titlerule[0.6pt]}]
\titlespacing*{\section}{0pt}{12pt}{6pt}

\newenvironment{entries}{\tabularx{\textwidth}{@{}p{2.3cm}X@{}}}{\endtabularx}
\setlist[itemize]{leftmargin=1.2em,itemsep=1pt,topsep=2pt}

\begin{document}

{\color{navy}\LARGE\bfseries Julio César Galindo López}\\[3pt]
Research Professor, Facultad de Empresariales, Universidad Panamericana, Ciudad UP, Mexico City\\
\href{mailto:jcgalindo@up.edu.mx}{jcgalindo@up.edu.mx} \;|\;
\href{https://ivlicaesar.github.io}{ivlicaesar.github.io} \;|\;
\href{https://github.com/IvliCaesar}{github.com/IvliCaesar}

\section{Research profile}
Research in pure mathematics at the interaction of algebraic geometry and homotopy theory: algebraic
cycles and motives (Peláez's orthogonal tower, levels of motives, the Bloch--Beilinson filtration),
rational homotopy theory and Galois actions, Sullivan's geometric topology, and frame functions,
Gleason's theorem and toposes. A second line applies random matrix theory, Fisher information and
optimal design to finance, genomics and imaging, together with applied econometric work on health
and transport economics. More than fifteen years of university teaching at UNAM, Universidad
Panamericana, Tecnológico de Monterrey and BUAP.

\section{Education}
\begin{entries}
2016--2021 & \textbf{Ph.D. in Mathematical Sciences}, Universidad Nacional Autónoma de México (UNAM), Instituto de Matemáticas.
Dissertation: \emph{Sobre la convergencia de la sucesión espectral ortogonal}. Advisor: Pablo Peláez Menaldo (Chebyshev Laboratory, St.~Petersburg State University).\\[3pt]
2014--2016 & \textbf{M.Sc. in Mathematical Sciences}, UNAM, Instituto de Matemáticas, Cuernavaca.
Thesis: \emph{La ley de reciprocidad de Drinfeld}. Advisor: Timothy Mooney Gendron Thornton.\\[3pt]
2008--2013 & \textbf{B.S. in Actuarial Science}, UNAM, Facultad de Ciencias.
Thesis: \emph{Random Matrices and $L$-Functions}. Advisor: Eduardo Dueñez Guzmán; co-advisor: Víctor Manuel Pérez-Abreu Carrión.\\[3pt]
2022 & \textbf{MITx MicroMasters Program in Statistics and Data Science}, Massachusetts Institute of Technology: \emph{Machine Learning with Python: From Linear Models to Deep Learning}; \emph{Probability: The Science of Uncertainty and Data}.\\[3pt]
2023--2024 & UNAM Diploma Program in Statistics and Data Mining, Module II: Probability and Stochastic Simulation.
\end{entries}

\section{Academic appointments}
\begin{entries}
Present & \textbf{Full-Time Research Professor}, Facultad de Empresariales, Universidad Panamericana, Ciudad UP. Courses: Fundamentos de Econometría, Estadística Inferencial, Cálculo Integral, Ecuaciones Diferenciales, Investigación de Operaciones, Cálculo Actuarial II, Probability and Statistics for Business.\\[3pt]
2023--2024 & Research Professor (Associate C), Benemérita Universidad Autónoma de Puebla, Facultad de Ciencias Físico Matemáticas: financial risk, portfolio theory and Monte Carlo methods.\\[3pt]
2023--present & Lecturer, IIMAS--FES Acatlán, UNAM: Statistical Computing; Probability and Stochastic Simulation.\\[3pt]
2016--present & Lecturer, Applied Mathematics and Computing, FES Acatlán, UNAM: Probability, Calculus III--IV, Logic, Geometry, Numerical Methods II.\\[3pt]
2021--2023 & Adjunct Professor, Tecnológico de Monterrey, Campus Ciudad de México: Software Development, Data Analytics and AI Tools, Data Science, Data Visualization, Natural Language Processing.\\[3pt]
2011--2022 & Lecturer, Facultad de Ciencias, UNAM: Probability I--II, Stochastic Processes I--II, Mathematical Analysis I, Measure Theory, Algebra, Calculus IV; seminars in actuarial applications and arithmetic geometry.\\[3pt]
2016--2017 & Teaching Assistant, Universidad Autónoma Metropolitana: Integral Calculus, Linear Algebra.
\end{entries}

\section{Publications}
\begin{itemize}
\item J.~C. Galindo and P. Peláez, \emph{On the convergence of the orthogonal spectral sequence and applications}, Homology, Homotopy and Applications \textbf{24} (2022), no.~2, 389--401. arXiv:2111.03508.
\item J.~C. Galindo, \emph{Convergence of the birational tower spectral sequence for mixed motives}, Homology, Homotopy and Applications.
\item J.~C. Galindo, \emph{La ley de reciprocidad de Drinfeld}, Morfismos (Cinvestav) \textbf{20} (2016), no.~2, 1--66.
\end{itemize}

\section{Submitted and in preparation}
\textit{Algebraic cycles, motives and homotopy theory}
\begin{itemize}
\item \emph{Levels of motives of smooth varieties}.
\item \emph{The orthogonal filtration on cycles of abelian varieties}.
\item \emph{A uniform orthogonal filtration on higher Chow groups}.
\item \emph{Divergence of the orthogonal tower of algebraic $K$-theory}.
\item \emph{On the Betti realization for the birational tower and applications}.
\item \emph{The orthogonal tower and the $E_1$-term of Voevodsky's slice spectral sequence}.
\item \emph{Weights, Koszul duality and Galois actions on rational homotopy}.
\item \emph{Motivic formality and non-abelian Galois actions on rational homotopy types}.
\item \emph{Formality loci in families and torsion of normal functions}.
\end{itemize}
\textit{Logic, toposes and geometry}
\begin{itemize}
\item with E. Dueñez and J. Iovino, \emph{Frame functions as deep computations: Gleason's theorem, tameness, and hidden variables}.
\item \emph{Frame functions on Riemannian manifolds}.
\item \emph{Homotopical and homological shadows of deep computation: solenoids, Frobenius, and persistent homology}.
\end{itemize}
\textit{Probability, statistics and applications}
\begin{itemize}
\item \emph{D-optimal design under additive Fisher information and the depth--panel trade-off for count data}.
\item \emph{One market, many eigenvalues: random matrix theory for equity panels with listings and delistings}.
\item with E. Hernández-Lemus, \emph{Which genes should a breast-cancer panel measure?}
\item with E. Hernández-Lemus, \emph{How much structure does a Hi-C map resolve? A Poisson random-matrix null, wavelet scales and D-optimal capture in triple-negative breast cancer}.
\item with L. Jiménez Casillas, \emph{Covariance estimation under portfolio constraints in Mexican equities}.
\item with M.~R. García Barrera, \emph{Spectral cleaning of high-dimensional covariance matrices for portfolio optimization: theory and cross-market evidence from Mexican and U.S. equities}.
\item with A. Kamei, \emph{Who sets the standard? Reference groups, medical information and weight misperception in Mexico}.
\item with X. Pacheco Zamora, \emph{Toll evasion, electronic payment and the economics of enforcement}.
\item with Ó. Amador, \emph{Which channels should an RGB camera add?}
\item \emph{One field, five readings: metaballs, Haar-wavelet texture, and the color grammar of a generative creature}, preprint.
\item with G. Villafuerte C., \emph{SIAMDA: a wavelet-based sentiment index for academic early warning}.
\item with M. Velázquez-Salazar, \emph{Reconfiguration of global value chains in Mexican coffee: labor, ethics, and sustainability from a humanistic perspective}.
\end{itemize}
\textit{Lecture notes}
\begin{itemize}
\item \emph{Localization, periodicity and Galois symmetry: an account of Sullivan's geometric topology}.
\item \emph{Three problems from Sullivan's geometric topology}.
\item \emph{Arithmetic homotopy theory and the Adams conjecture: notes on Sullivan's Genetics}.
\item \emph{Gaussian random matrix ensembles: classical theory and universality}, after a course of E. Dueñez at CIMAT.
\item \emph{Wavelet theory}, lecture notes.
\end{itemize}

\section{Talks}
\begin{entries}
2026 & \emph{La IA como aliada del emprendedor}, XXVII Semana de la Facultad de Economía, Universidad Autónoma del Estado de México.\\
2026 & \emph{¿Qué matemáticas necesitas para el Machine Learning?}, Facultad de Empresariales, Universidad Panamericana.\\
2026 & \emph{Homología persistente aplicada a datos actuariales}, CONAC, Universidad Panamericana.\\
2026 & \emph{De bordes a leyes de grupo: cobordismo complejo}, colloquium, AGATA, Cuernavaca.\\
2026 & \emph{Reconfiguration of global value chains in Mexican coffee} (with M. Velázquez-Salazar), WSSA.\\
2019 & 52nd Congress of the Mexican Mathematical Society, Universidad Autónoma de Nuevo León.\\
2018 & \emph{Ciclos algebraicos y variedades abelianas}, Seminario de Becarios, Instituto de Matemáticas, UNAM.\\
2018 & Motives in St.~Petersburg, Euler International Mathematical Institute, St.~Petersburg.\\
2017 & \emph{Matrices aleatorias y teoría de números}, Coloquio de Orientación Matemática, Facultad de Ciencias, UNAM.\\
2016 & \emph{Teoría de martingalas a tiempo discreto}, Seminario Movimiento Browniano, Martingalas y Cálculo Estocástico, UNAM.\\
2016 & CIMPA--CIMAT--ICTP Summer School on Moduli of Curves, CIMAT; National Congress of Algebraic Geometry, Casa Matemática Oaxaca.\\
2014 & Applications of Étale Cohomology Summer School, Alfréd Rényi Institute of Mathematics, Budapest.\\
2011 & Mini-course \emph{Random Matrices and Number Theory}, Third Interinstitutional Seminar on Random Matrices, CIMAT; 35th Conference on Stochastic Processes and their Applications, Oaxaca.\\
2010 & \emph{Random Matrices and the Riemann Zeta Function}, Second Interinstitutional Seminar on Random Matrices, CIMAT.
\end{entries}

\section{Student supervision}
\begin{entries}
2026--present & Co-advisor (with L. Jiménez Casillas), Daniela Cojab Mizrahi, B.A., Universidad Panamericana: AI, homogenization and fragility in equity markets.\\
2026--present & Reader, Emma Carolina Alfaro Ramírez, doctoral protocol, Doctorado en Ciencias Empresariales, Universidad Panamericana.\\
2027 (expected) & Advisor, Ximena Pacheco Zamora, B.Sc., Facultad de Ciencias, UNAM: electronic-toll adoption and payment evasion.\\
2025--2026 & Advisor, Martín Rodolfo García Barrera, B.Sc., UNAM: covariance estimation via random matrix theory and portfolio optimization.\\
2025 & Reader, Miguel Ángel Rodríguez Islas, M.Sc., UNAM: \emph{Clases de Chern en geometría algebraica}.\\
2024 & Advisor, José Alberto Meléndez Piña, M.Sc., IMATE--UNAM: the Fourier--Mukai transform.\\
2024 & Advisor, Francisco Estrada Ruano, B.Sc., FES Acatlán, UNAM: random matrices and determinantal processes.\\
2023 & Advisor, Fernando Flores Tiburcio, B.Sc., Facultad de Ciencias, UNAM: almost coherent sheaves on compact Riemann surfaces (honorable mention).\\
2014, 2024 & Committee member: Eduardo Granados González (random-number generator from an air lottery machine); Martín Guarneros Martínez ($p$-adic differential calculus).
\end{entries}

\section{Collaborations and consulting}
\begin{entries}
2026 & SIAMDA, with G. Villafuerte C. (FES Acatlán, UNAM); AI adoption and organizational KPIs, with S. García Álvarez (Universidad Panamericana).\\
2023--2024 & Carl Zeiss Group: insurance portfolio optimization, pricing and risk mitigation.\\
2023--2024 & Comisión Nacional de Seguros y Fianzas: premium forecasting with time-series methods; sentiment analysis and boosting for gender-related policy analysis.\\
2023--2024 & Review of the Calculus I and IV curricula, Applied Mathematics and Computing, FES Acatlán.
\end{entries}

\section{Honors and service}
\begin{entries}
2024--present & Candidate, Sistema Nacional de Investigadores (SNI).\\
2024 & Recognition for academic service, FES Acatlán; UNAM recognition for ten years of academic service.\\
2023 & Referee, \emph{Advances in Mathematics}. Working group for the M.Sc. in Data Science Research, BUAP.\\
2019--2023 & Doctoral scholarship, CONACyT.\\
2016--2019 & Master's scholarship, CONACyT.\\
2009--2011 & Research assistant (SNI III), CIMAT.
\end{entries}

\section{Languages and computing}
Spanish (native); English, French and Russian (intermediate).
R, Python, Julia, SQL, C, C++, Java, Matlab, Mathematica, Lean, \LaTeX, HTML.

\end{document}
